\section{Introduction}

This document develops an improved electoral procedure for the Bundestag. Although
the proposed procedure is primarily intended for the election of the German Bundestag, it can
be applied to a wide variety of other elections. It is essentially a general
electoral procedure developed using the Bundestag election as an example.

There are several points of criticism regarding the current electoral law:
\begin{enumerate}
  \item Forced grudging coalitions with shabby compromises and many disputes.
    Given the fragmented party landscape, it is hardly possible for a single party to win an absolute
    majority, so coalitions must be formed. However, coalition partners often block each other
    to such an extent that nothing gets accomplished and all partners end up looking like liars.

  \item Lead candidates are not elected by the citizens. Those who vote for a particular party based on
    its positions must also accept its lead candidate --- regardless of who that person is.

  \item Under the current implementation of personalized proportional representation, it can
    happen that some constituencies are not personally represented in the Bundestag. This was
    criticized by Friedrich Merz after the 2025 federal election. The previous system, on the other hand,
    had the widely criticized disadvantage of an ever-expanding Bundestag.
\end{enumerate}

All three points of criticism are avoided by the system derived in the following.

The proposal should be understood as a contribution to a discussion that still needs to be initiated
across society, which could end with a constitutional reform or a new constitution.

To support the description of the procedure, there is a simulation available as a public
\href{https://github.com/danielschwalbe/ipres}{GitHub project.} The included Jupyter notebooks
can be used to experiment with the procedure and study the effects
of various parameters.

The simulation is also the ultimate procedural reference, as it is tested and furthermore
has a level of detail that is very difficult to achieve in such a Google Docs document. This document, on the other hand, explains the why and wherefore. This is only included in a very
condensed form in the simulation documentation.
